\section{Affine schemes}
\label{section:I.1}

\setcounter{subsection}{0}
\subsection{The prime spectrum of a ring}
\label{subsection:I.1.1}

\oldpage[I]{80}
\begin{env}[1.1.1]
\label{I.1.1.1}
\emph{Notation}: Let $A$ be a (commutative) ring and $M$ an $A$-module. In
this chapter and the following ones, we shall constantly use the following
notation:

\smallskip
\noindent
$\operatorname{Spec}(A)={}$\emph{set of prime ideals} of $A$, also called
the \emph{prime spectrum} of $A$; for an
$x\in X=\operatorname{Spec}(A)$, it will often be convenient to write
$\mathfrak{j}_x$ instead of $x$. For $\operatorname{Spec}(A)$ to be
\emph{empty}, it is necessary and sufficient that the ring $A$ consist of
$0$ alone.

\smallskip
\noindent
$A_x=A_{\mathfrak{j}_x}={}$\emph{(local) ring of fractions $S^{-1}A$}, where
$S=A-\mathfrak{j}_x$.

\smallskip
\noindent
$\mathfrak{m}_x=\mathfrak{j}_xA_{\mathfrak{j}_x}={}$\emph{maximal ideal of
$A_x$}.

\smallskip
\noindent
$k(x)=A_x/\mathfrak{m}_x={}$\emph{residue field of $A_x$}, canonically
isomorphic to the field of fractions of the integral domain
$A/\mathfrak{j}_x$, with which it is identified.

\smallskip
\noindent
$f(x)={}$\emph{class of $f$ mod.~$\mathfrak{j}_x$}, in
$A/\mathfrak{j}_x\subset k(x)$, for $f\in A$ and $x\in X$. We also say that
$f(x)$ is the \emph{value} of $f$ at the point $x\in\operatorname{Spec}(A)$;
the relations $f(x)=0$ and $f\in\mathfrak{j}_x$ are \emph{equivalent}.

\smallskip
\noindent
$M_x=M\otimes_AA_x={}$\emph{module of fractions with denominators in
$A-\mathfrak{j}_x$}.

\smallskip
\noindent
$\mathfrak{r}(E)={}$\emph{radical of the ideal of $A$ generated by a subset
$E$ of $A$}.

\smallskip
\noindent
$V(E)={}$\emph{set of $x\in X$ such that $E\subset\mathfrak{j}_x$}
(or equivalently the set of $x\in X$ such that $f(x)=0$ for all $f\in E$),
for $E\subset A$. We therefore have
\begin{equation}
  \mathfrak{r}(E)=\bigcap_{x\in V(E)}\mathfrak{j}_x.
  \tag{1.1.1.1}\label{I.1.1.1.1}
\end{equation}

\smallskip
\noindent
$V(f)=V(\{f\})$ for $f\in A$.

\smallskip
\noindent
$D(f)=X-V(f)={}$\emph{set of $x\in X$ where $f(x)\neq0$}.
\end{env}

\begin{proposition}[1.1.2]
\label{I.1.1.2}
We have the following properties:
\begin{enumerate}
  \item[\rm(i)] $V(0)=X$, $V(1)=\varnothing$.
  \item[\rm(ii)] The relation $E\subset E'$ implies $V(E)\supset V(E')$.
  \item[\rm(iii)] For every family $(E_\lambda)$ of subsets of $A$,
  \[
    V\left(\bigcup_\lambda E_\lambda\right)
      =V\left(\sum_\lambda E_\lambda\right)
      =\bigcap_\lambda V(E_\lambda).
  \]
  \item[\rm(iv)] $V(EE')=V(E)\cup V(E')$.
  \item[\rm(v)] $V(E)=V(\mathfrak{r}(E))$.
\end{enumerate}
\end{proposition}

Properties (i), (ii), (iii) are trivial, and (v) follows from (ii) and from
formula (\hyperref[I.1.1.1.1]{1.1.1.1}). It is clear that
$V(EE')\supset V(E)\cup V(E')$; conversely, if $x\notin V(E)$ and
$x\notin V(E')$, there exist $f\in E$ and $f'\in E'$ such that $f(x)\neq0$
and $f'(x)\neq0$ in $k(x)$, whence $f(x)f'(x)\neq0$, in other words
$x\notin V(EE')$, which proves (iv).

Proposition (\hyperref[I.1.1.2]{1.1.2}) shows, among other things, that
the sets of the form $V(E)$ (where $E$ runs through the set of subsets of $A$)
are the \emph{closed sets} of a topology on $X$, which we shall call the
\emph{spectral topology}\footnote{The introduction of this topology into
algebraic geometry is due to Zariski. It is therefore often called the
``Zariski topology'' of $X$.}; unless expressly stated otherwise, we shall
always assume $X=\operatorname{Spec}(A)$ endowed with the spectral topology.

\oldpage[I]{81}
\begin{env}[1.1.3]
\label{I.1.1.3}
For every subset $Y$ of $X$, we shall denote by $\mathfrak{j}(Y)$
the set of $f\in A$ such that $f(y)=0$ for all $y\in Y$; it amounts to the
same to say that $\mathfrak{j}(Y)$ is the intersection of the prime ideals
$\mathfrak{j}_y$ for $y\in Y$. It is clear that the relation $Y\subset Y'$
implies $\mathfrak{j}(Y)\supset\mathfrak{j}(Y')$ and that we have
\begin{equation}
  \mathfrak{j}\left(\bigcup_\lambda Y_\lambda\right)
    =\bigcap_\lambda\mathfrak{j}(Y_\lambda)
  \tag{1.1.3.1}\label{I.1.1.3.1}
\end{equation}
for every family $(Y_\lambda)$ of subsets of $X$. Finally, we have
\begin{equation}
  \mathfrak{j}(\{x\})=\mathfrak{j}_x.
  \tag{1.1.3.2}\label{I.1.1.3.2}
\end{equation}
\end{env}

\begin{proposition}[1.1.4]
\label{I.1.1.4}
\begin{enumerate}
  \item[\rm(i)] For every subset $E$ of $A$, we have
  $\mathfrak{j}(V(E))=\mathfrak{r}(E)$.
  \item[\rm(ii)] For every subset $Y$ of $X$,
  $V(\mathfrak{j}(Y))=\overline{Y}$, the closure of $Y$ in $X$.
\end{enumerate}
\end{proposition}

(i) is an immediate consequence of the definitions and of
(\hyperref[I.1.1.1.1]{1.1.1.1}); on the other hand,
$V(\mathfrak{j}(Y))$ is closed and contains $Y$; conversely, if
$Y\subset V(E)$, we have $f(y)=0$ for all $f\in E$ and all $y\in Y$, hence
$E\subset\mathfrak{j}(Y)$, $V(E)\supset V(\mathfrak{j}(Y))$, which proves
(ii).

\begin{corollary}[1.1.5]
\label{I.1.1.5}
The closed subsets of $X=\operatorname{Spec}(A)$ and the ideals of $A$ equal
to their radicals (in other words, the intersections of prime ideals)
correspond one-to-one under the decreasing maps
$Y\mapsto\mathfrak{j}(Y)$, $\mathfrak{a}\mapsto V(\mathfrak{a})$; to the
union $Y_1\cup Y_2$ of two closed subsets corresponds
$\mathfrak{j}(Y_1)\cap\mathfrak{j}(Y_2)$, and to the intersection of an
arbitrary family $(Y_\lambda)$ of closed subsets corresponds the radical of
the sum of the $\mathfrak{j}(Y_\lambda)$.
\end{corollary}

\begin{corollary}[1.1.6]
\label{I.1.1.6}
If $A$ is a Noetherian ring, $X=\operatorname{Spec}(A)$ is a Noetherian
space.
\end{corollary}

Note that the converse of this corollary is false, as is shown by the
example of a non-Noetherian integral domain having only one prime ideal
$\neq\{0\}$, for example a non-discrete valuation ring of rank 1.

As an example of a ring $A$ whose spectrum is not a Noetherian space, one
may mention the ring $\mathcal{C}(Y)$ of continuous real-valued functions on
an infinite compact space $Y$; it is known that, as a set, $Y$ can be
identified with the set of maximal ideals of $A$, and it is easy to see that
the topology induced on $Y$ by that of $X=\operatorname{Spec}(A)$ is the
original topology of $Y$. Since $Y$ is not a Noetherian space, the same is
therefore true of $X$.

\begin{corollary}[1.1.7]
\label{I.1.1.7}
For every $x\in X$, the closure of $\{x\}$ is the set of $y\in X$ such
that $\mathfrak{j}_x\subset\mathfrak{j}_y$. For $\{x\}$ to be closed, it is
necessary and sufficient that $\mathfrak{j}_x$ be maximal.
\end{corollary}

\begin{corollary}[1.1.8]
\label{I.1.1.8}
The space $X=\operatorname{Spec}(A)$ is a Kolmogorov space.
\end{corollary}

Indeed, if $x$, $y$ are two distinct points of $X$, we have either
$\mathfrak{j}_x\not\subset\mathfrak{j}_y$ or
$\mathfrak{j}_y\not\subset\mathfrak{j}_x$, hence one of the points $x$, $y$
does not belong to the closure of the other.

\begin{env}[1.1.9]
\label{I.1.1.9}
By Proposition (\hyperref[I.1.1.2]{1.1.2}, (iv)), for two elements
$f$, $g$ of $A$, we have
\begin{equation}
  D(fg)=D(f)\cap D(g).
  \tag{1.1.9.1}\label{I.1.1.9.1}
\end{equation}
Note also that the relation $D(f)=D(g)$ means, by Proposition
(\hyperref[I.1.1.4]{1.1.4}, (i)) and Proposition
(\hyperref[I.1.1.2]{1.1.2}, (v)), that
$\mathfrak{r}(f)=\mathfrak{r}(g)$, or equivalently that the minimal prime
ideals containing $(f)$ and $(g)$ are the same; in particular, this is the
case when $f=ug$, where $u$ is invertible.
\end{env}

\oldpage[I]{82}
\begin{proposition}[1.1.10]
\label{I.1.1.10}
\begin{enumerate}
  \item[\rm(i)] As $f$ runs through $A$, the sets $D(f)$ form a
  base for the topology of $X$.
  \item[\rm(ii)] For every $f\in A$, $D(f)$ is quasi-compact. In particular,
  $X=D(1)$ is quasi-compact.
\end{enumerate}
\end{proposition}

(i) Let $U$ be an open set in $X$; by definition, we have
$U=X-V(E)$, where $E$ is a subset of $A$, and
$V(E)=\bigcap_{f\in E}V(f)$, whence $U=\bigcup_{f\in E}D(f)$.

(ii) By (i), it suffices to prove that if $(f_\lambda)_{\lambda\in L}$
is a family of elements of $A$ such that
$D(f)\subset\bigcup_{\lambda\in L}D(f_\lambda)$, there exists a finite subset
$J$ of $L$ such that $D(f)\subset\bigcup_{\lambda\in J}D(f_\lambda)$. Let
$\mathfrak{a}$ be the ideal of $A$ generated by the $f_\lambda$; by
hypothesis we have $V(f)\supset V(\mathfrak{a})$, hence
$\mathfrak{r}(f)\subset\mathfrak{r}(\mathfrak{a})$; since
$f\in\mathfrak{r}(f)$, there exists an integer $n\geq0$ such that
$f^n\in\mathfrak{a}$. But then $f^n$ belongs to the ideal $\mathfrak{b}$
generated by a finite subfamily $(f_\lambda)_{\lambda\in J}$, and we have
$V(f)=V(f^n)\supset V(\mathfrak{b})=\bigcap_{\lambda\in J}V(f_\lambda)$,
that is to say, $D(f)\subset\bigcup_{\lambda\in J}D(f_\lambda)$.

\begin{proposition}[1.1.11]
\label{I.1.1.11}
For every ideal $\mathfrak{a}$ of $A$,
$\operatorname{Spec}(A/\mathfrak{a})$ is canonically identified with the
closed subspace $V(\mathfrak{a})$ of $\operatorname{Spec}(A)$.
\end{proposition}

Indeed, it is known that there is a canonical one-to-one correspondence,
respecting the order structure given by inclusion, between the ideals (resp.
prime ideals) of $A/\mathfrak{a}$ and the ideals (resp. prime ideals) of $A$
containing $\mathfrak{a}$.

Recall that the set $\mathfrak{N}$ of nilpotent elements of $A$ (or
\emph{nilradical} of $A$) is an ideal equal to $\mathfrak{r}(0)$, the
intersection of all the prime ideals of $A$
(\hyperref[0.1.1.1]{0, 1.1.1}).

\begin{corollary}[1.1.12]
\label{I.1.1.12}
The topological spaces $\operatorname{Spec}(A)$ and
$\operatorname{Spec}(A/\mathfrak{N})$ are canonically isomorphic.
\end{corollary}

\begin{proposition}[1.1.13]
\label{I.1.1.13}
For $X=\operatorname{Spec}(A)$ to be irreducible
(\hyperref[0.2.1.1]{0, 2.1.1}), it is necessary and sufficient that the
ring $A/\mathfrak{N}$ be an integral domain (or, what amounts to the same,
that the ideal $\mathfrak{N}$ be prime).
\end{proposition}

By Corollary (\hyperref[I.1.1.12]{1.1.12}), we may restrict ourselves to
the case where $\mathfrak{N}=0$. If $X$ is reducible, there exist two closed
subsets $Y_1$, $Y_2$ distinct from $X$ such that $X=Y_1\cup Y_2$, whence
$\mathfrak{j}(X)=\mathfrak{j}(Y_1)\cap\mathfrak{j}(Y_2)=0$, the ideals
$\mathfrak{j}(Y_1)$ and $\mathfrak{j}(Y_2)$ being distinct from $(0)$
(\hyperref[I.1.1.5]{1.1.5}); hence $A$ is not an integral domain.
Conversely, if there are in $A$ two elements $f\neq0$, $g\neq0$ such that
$fg=0$, we have $V(f)\neq X$, $V(g)\neq X$ (since the intersection of the
prime ideals of $A$ is $(0)$), and $X=V(fg)=V(f)\cup V(g)$.

\begin{corollary}[1.1.14]
\label{I.1.1.14}
\begin{enumerate}
  \item[\rm(i)] Under the one-to-one correspondence between closed subsets
  of $X=\operatorname{Spec}(A)$ and ideals of $A$ equal to their radicals,
  the irreducible closed subsets of $X$ correspond to the prime ideals of
  $A$. In particular, the irreducible components of $X$ correspond to the
  minimal prime ideals of $A$.
  \item[\rm(ii)] The map $x\mapsto\overline{\{x\}}$ establishes a
  one-to-one correspondence between $X$ and the set of irreducible closed
  subsets of $X$ (in other words, every irreducible closed subset of $X$
  has one and only one generic point).
\end{enumerate}
\end{corollary}

(i) follows at once from (\hyperref[I.1.1.13]{1.1.13}) and
(\hyperref[I.1.1.11]{1.1.11}); and to prove (ii) we may, by virtue of
(\hyperref[I.1.1.11]{1.1.11}), restrict ourselves to the case where $X$ is
irreducible; then, by Proposition (\hyperref[I.1.1.13]{1.1.13}), there
exists in $A$ a smallest prime ideal $\mathfrak{N}$, which therefore
corresponds to a generic point
\oldpage[I]{83}
of $X$; moreover, $X$ has only one generic point, since it is a
Kolmogorov space ((\hyperref[I.1.1.8]{1.1.8}) and
(\hyperref[0.2.1.3]{0, 2.1.3})).

\begin{proposition}[1.1.15]
\label{I.1.1.15}
If $\mathfrak{J}$ is an ideal of $A$ contained in the radical
$\mathfrak{R}(A)$, the only neighbourhood of $V(\mathfrak{J})$ in
$X=\operatorname{Spec}(A)$ is the whole space $X$.
\end{proposition}

Indeed, every maximal ideal of $A$ belongs by definition to
$V(\mathfrak{J})$. Since every ideal $\mathfrak{a}\neq A$ of $A$ is contained
in a maximal ideal, we have $V(\mathfrak{a})\cap V(\mathfrak{J})\neq\varnothing$,
whence the proposition.

\subsection{Functorial properties of the prime spectra of rings}
\label{subsection:I.1.2}

\begin{env}[1.2.1]
\label{I.1.2.1}
Let $A$, $A'$ be two rings and
\[
  \varphi:A'\longrightarrow A
\]
a ring homomorphism. For every prime ideal
$x=\mathfrak{j}_x\in\operatorname{Spec}(A)=X$, the ring
$A'/\varphi^{-1}(\mathfrak{j}_x)$ is canonically isomorphic to a
subring of $A/\mathfrak{j}_x$, hence is an integral domain; in other words,
$\varphi^{-1}(\mathfrak{j}_x)$ is a prime ideal of $A'$; we shall
denote it by ${}^a\varphi(x)$, and we have thus defined a map
\[
  {}^a\varphi:X=\operatorname{Spec}(A)\longrightarrow
  X'=\operatorname{Spec}(A')
\]
(also denoted $\operatorname{Spec}(\varphi)$), which we shall call the map
associated with the homomorphism $\varphi$. We shall denote by
$\varphi^x$ the injective homomorphism of
$A'/\varphi^{-1}(\mathfrak{j}_x)$ into $A/\mathfrak{j}_x$ deduced from
$\varphi$ by passing to the quotients, as well as its canonical extension to
a monomorphism of fields
\[
  \varphi^x:k({}^a\varphi(x))\longrightarrow k(x)~;
\]
for every $f'\in A'$, we therefore have by definition
\begin{equation}
  \varphi^x\bigl(f'({}^a\varphi(x))\bigr)=\bigl(\varphi(f')\bigr)(x)
  \qquad(x\in X).
  \tag{1.2.1.1}\label{I.1.2.1.1}
\end{equation}
\end{env}

\begin{proposition}[1.2.2]
\label{I.1.2.2}
\begin{enumerate}
  \item[\rm(i)] For every subset $E'$ of $A'$, we have
  \begin{equation}
    {}^a\varphi^{-1}(V(E'))=V(\varphi(E'))
    \tag{1.2.2.1}\label{I.1.2.2.1}
  \end{equation}
  and in particular, for every $f'\in A'$,
  \begin{equation}
    {}^a\varphi^{-1}(D(f'))=D(\varphi(f')).
    \tag{1.2.2.2}\label{I.1.2.2.2}
  \end{equation}
  \item[\rm(ii)] For every ideal $\mathfrak{a}$ of $A$, we have
  \begin{equation}
    \overline{{}^a\varphi(V(\mathfrak{a}))}
    =V(\varphi^{-1}(\mathfrak{a})).
    \tag{1.2.2.3}\label{I.1.2.2.3}
  \end{equation}
\end{enumerate}
\end{proposition}

Indeed, the relation ${}^a\varphi(x)\in V(E')$ is by definition
equivalent to $E'\subset\varphi^{-1}(\mathfrak{j}_x)$, hence to
$\varphi(E')\subset\mathfrak{j}_x$, and finally to
$x\in V(\varphi(E'))$, whence (i). To prove (ii), we may assume
$\mathfrak{a}$ equal to its radical, since
$V(\mathfrak{r}(\mathfrak{a}))=V(\mathfrak{a})$
(\hyperref[I.1.1.2]{1.1.2 (v)}) and
$\varphi^{-1}(\mathfrak{r}(\mathfrak{a}))
=\mathfrak{r}(\varphi^{-1}(\mathfrak{a}))$; if we set
$Y=V(\mathfrak{a})$ and
$\mathfrak{a}'=\mathfrak{j}({}^a\varphi(Y))$, we have
$\overline{{}^a\varphi(Y)}=V(\mathfrak{a}')$ (Proposition
(\hyperref[I.1.1.4]{1.1.4 (ii)})); the relation
$f'\in\mathfrak{a}'$ is by definition equivalent to $f'(x')=0$ for all
$x'\in{}^a\varphi(Y)$; hence, by formula
(\hyperref[I.1.2.1.1]{1.2.1.1}), it is also equivalent to
$\varphi(f')(x)=0$ for all $x\in Y$, or again to
$\varphi(f')\in\mathfrak{j}(Y)=\mathfrak{a}$, since $\mathfrak{a}$ is
equal to its radical; whence (ii).

\begin{corollary}[1.2.3]
\label{I.1.2.3}
The map ${}^a\varphi$ is continuous.
\end{corollary}

Observe that if $A''$ is a third ring and $\varphi'$ a homomorphism
$A''\to A'$, we have
${}^a(\varphi\circ\varphi')={}^a\varphi'\circ{}^a\varphi$; this result and
Corollary (\hyperref[I.1.2.3]{1.2.3}) mean that
$\operatorname{Spec}(A)$ is a contravariant functor in $A$, from the
category of rings to that of topological spaces.

\begin{corollary}[1.2.4]
\label{I.1.2.4}
\oldpage[I]{84}
Suppose that $\varphi$ is such that every $f\in A$ can be written
$f=h\varphi(f')$, where $h$ is invertible in $A$ (which is in particular
the case when $\varphi$ is surjective). Then ${}^a\varphi$ is a
homeomorphism of $X$ onto ${}^a\varphi(X)$.
\end{corollary}

Let us show that for every subset $E\subset A$, there exists a subset $E'$
of $A'$ such that $V(E)=V(\varphi(E'))$; by virtue of axiom $(T_0)$
(\hyperref[I.1.1.8]{1.1.8}) and of formula
(\hyperref[I.1.2.2.1]{1.2.2.1}), this will imply first that
${}^a\varphi$ is injective, and then, again by
(\hyperref[I.1.2.2.1]{1.2.2.1}), that ${}^a\varphi$ is a
homeomorphism. Now, for each $f\in E$ it suffices to take an $f'\in A'$
such that $h\varphi(f')=f$ with $h$ invertible in $A$; the set $E'$ of these
elements $f'$ has the required property.

\begin{env}[1.2.5]
\label{I.1.2.5}
In particular, when $\varphi$ is the canonical homomorphism of $A$ onto a
quotient ring $A/\mathfrak{a}$, we recover
(\hyperref[I.1.1.12]{1.1.12}), and ${}^a\varphi$ is the canonical
injection of $V(\mathfrak{a})$, identified with
$\operatorname{Spec}(A/\mathfrak{a})$, into $X=\operatorname{Spec}(A)$.
\end{env}

Another special case of (\hyperref[I.1.2.4]{1.2.4}) gives:

\begin{corollary}[1.2.6]
\label{I.1.2.6}
If $S$ is a multiplicative subset of $A$, the spectrum
$\operatorname{Spec}(S^{-1}A)$ is canonically identified (together with its
topology) with the subspace of $X=\operatorname{Spec}(A)$ consisting of the
$x$ such that $\mathfrak{j}_x\cap S=\varnothing$.
\end{corollary}

Indeed, it is known (\hyperref[0.1.2.6]{0, 1.2.6}) that the prime ideals
of $S^{-1}A$ are the ideals $S^{-1}\mathfrak{j}_x$ such that
$\mathfrak{j}_x\cap S=\varnothing$, and that we have
$\mathfrak{j}_x=({}^S i_A)^{-1}(S^{-1}\mathfrak{j}_x)$. It therefore
suffices to apply Corollary (\hyperref[I.1.2.4]{1.2.4}) to ${}^S i_A$.

\begin{corollary}[1.2.7]
\label{I.1.2.7}
For ${}^a\varphi(X)$ to be dense in $X'$, it is necessary and sufficient
that every element of the kernel $\operatorname{Ker}\varphi$ be nilpotent.
\end{corollary}

Indeed, applying formula (\hyperref[I.1.2.2.3]{1.2.2.3}) to the ideal
$\mathfrak{a}=(0)$, we get
$\overline{{}^a\varphi(X)}=V(\operatorname{Ker}\varphi)$, and for
$V(\operatorname{Ker}\varphi)=X'$ it is necessary and sufficient that
$\operatorname{Ker}\varphi$ be contained in all the prime ideals of
$A'$, that is to say, in the nilradical $\mathfrak{N}'$ of $A'$.

\subsection{Sheaf associated with a module}
\label{subsection:I.1.3}
\label{I.1.3}

\begin{env}[1.3.1]
\label{I.1.3.1}
Let $A$ be a commutative ring, $M$ an $A$-module, $f$ an element of $A$, and
$S_f$ the multiplicative set of the $f^n$, where $n\geq0$. Recall that we
set $A_f=S_f^{-1}A$, $M_f=S_f^{-1}M$. If $S'_f$ is the saturated
multiplicative subset of $A$ consisting of the $g\in A$ that divide an
element of $S_f$, it is known that $A_f$ and $M_f$ are canonically identified
with ${S'_f}^{-1}A$ and ${S'_f}^{-1}M$
(\hyperref[0.1.4.3]{0, 1.4.3}).
\end{env}

\begin{lemma}[1.3.2]
\label{I.1.3.2}
The following conditions are equivalent:
\begin{enumerate}
  \item[\rm(a)] $g\in S'_f$;
  \item[\rm(b)] $S'_g\subset S'_f$;
  \item[\rm(c)] $f\in\mathfrak{r}(g)$;
  \item[\rm(d)] $\mathfrak{r}(f)\subset\mathfrak{r}(g)$;
  \item[\rm(e)] $V(g)\subset V(f)$;
  \item[\rm(f)] $D(f)\subset D(g)$.
\end{enumerate}
\end{lemma}

This follows immediately from the definitions and from
(\hyperref[I.1.1.5]{1.1.5}).

\begin{env}[1.3.3]
\label{I.1.3.3}
If $D(f)=D(g)$, Lemma (\hyperref[I.1.3.2]{1.3.2, b)}) shows that
$M_f=M_g$. More generally, if $D(f)\supset D(g)$, and hence
$S'_f\subset S'_g$, it is known (\hyperref[0.1.4.1]{0, 1.4.1}) that there
exists a canonical functorial homomorphism
\[
  \rho_{g,f}:M_f\longrightarrow M_g,
\]
and if $D(f)\supset D(g)\supset D(h)$, we have
(\hyperref[0.1.4.4]{0, 1.4.4})
\begin{equation}
  \rho_{h,g}\circ\rho_{g,f}=\rho_{h,f}.
  \tag{1.3.3.1}\label{I.1.3.3.1}
\end{equation}
\end{env}

\oldpage[I]{85}
As $f$ runs through $A-\mathfrak{j}_x$ (for a given $x$ in
$X=\operatorname{Spec}(A)$), the sets $S'_f$ form an increasing directed set
of subsets of $A-\mathfrak{j}_x$, since for two elements
$f$, $g$ of $A-\mathfrak{j}_x$, $S'_f$ and $S'_g$ are contained in
$S'_{fg}$; since the union of the $S'_f$ for $f\in A-\mathfrak{j}_x$ is
$A-\mathfrak{j}_x$, we conclude (\hyperref[0.1.4.5]{0, 1.4.5}) that the
$A_x$-module $M_x$ is canonically identified with the \emph{inductive limit}
$\varinjlim M_f$ with respect to the family of homomorphisms
$(\rho_{g,f})$. We shall denote by
\[
  \rho_x^f:M_f\longrightarrow M_x
\]
the canonical homomorphism for $f\in A-\mathfrak{j}_x$ (or, what amounts to
the same, $x\in D(f)$).

\begin{definition}[1.3.4]
\label{I.1.3.4}
We call \emph{structure sheaf of the prime spectrum}
$X=\operatorname{Spec}(A)$ (resp. \emph{sheaf associated with the $A$-module}
$M$), and denote by $\widetilde{A}$ or $\mathscr{O}_X$ (resp. $\widetilde{M}$),
the sheaf of rings (resp. the $\widetilde{A}$-module) associated with the
presheaf $D(f)\mapsto A_f$ (resp. $D(f)\mapsto M_f$) on the base $\mathfrak{B}$
of $X$ consisting of the $D(f)$, where $f\in A$
((\hyperref[I.1.1.10]{1.1.10}),
(\hyperref[0.3.2.1]{0, 3.2.1} and
\hyperref[0.3.5.6]{3.5.6})).
\end{definition}

We have seen (\hyperref[0.3.2.4]{0, 3.2.4}) that the stalk
$\widetilde{A}_x$ (resp. $\widetilde{M}_x$) \emph{is identified with the ring}
$A_x$ (resp. \emph{with the $A_x$-module} $M_x$); we shall denote by
\[
  \theta_f:A_f\longrightarrow\Gamma(D(f),\widetilde{A})
\]
\[
  \text{(resp. }\theta_f:M_f\longrightarrow
  \Gamma(D(f),\widetilde{M})\text{),}
\]
the canonical map, so that for every $x\in D(f)$ and every
$\xi\in M_f$, we have
\begin{equation}
  (\theta_f(\xi))_x=\rho_x^f(\xi).
  \tag{1.3.4.1}\label{I.1.3.4.1}
\end{equation}

\begin{proposition}[1.3.5]
\label{I.1.3.5}
$\widetilde{M}$ is an exact covariant functor in $M$, from the category of
$A$-modules to the category of $\widetilde{A}$-modules.
\end{proposition}

Indeed, let $M$, $N$ be two $A$-modules and $u$ a homomorphism
$M\to N$; for every $f\in A$, there corresponds canonically to $u$ a
homomorphism $u_f$ of the $A_f$-module $M_f$ into the $A_f$-module $N_f$,
and the diagram (for $D(g)\subset D(f)$)
\[
  \xymatrix{
    M_f\ar[r]^{u_f}\ar[d]_{\rho_{g,f}} & N_f\ar[d]^{\rho_{g,f}}\\
    M_g\ar[r]_{u_g} & N_g
  }
\]
is commutative (\hyperref[0.1.4.1]{0, 1.4.1}); these homomorphisms
therefore define a homomorphism of $\widetilde{A}$-modules
$\widetilde{u}:\widetilde{M}\to\widetilde{N}$
(\hyperref[0.3.2.3]{0, 3.2.3}). Moreover, for every $x\in X$,
$\widetilde{u}_x$ is the inductive limit of the $u_f$ for $x\in D(f)$
($f\in A$), and consequently (\hyperref[0.1.4.5]{0, 1.4.5}), if we
canonically identify $\widetilde{M}_x$ and $\widetilde{N}_x$ with $M_x$ and
$N_x$ respectively, $\widetilde{u}_x$ is identified with the homomorphism
$u_x$ deduced canonically from $u$. If $P$ is a third $A$-module, $v$ a
homomorphism $N\to P$ and $w=v\circ u$, it is immediate that
$w_x=v_x\circ u_x$, hence $\widetilde{w}=\widetilde{v}\circ\widetilde{u}$.
We have therefore indeed defined a \emph{covariant functor} $\widetilde{M}$
in $M$, from the category of $A$-modules to that of $\widetilde{A}$-modules.
\emph{This functor is exact}, since for every $x\in X$, $M_x$ is an exact
functor in $M$ (\hyperref[0.1.3.2]{0, 1.3.2}); moreover, we have
$\operatorname{Supp}(M)=\operatorname{Supp}(\widetilde{M})$ by the
definitions of the two sides (\hyperref[0.1.7.1]{0, 1.7.1} and
\hyperref[0.3.1.6]{3.1.6}).

\oldpage[I]{86}
\begin{proposition}[1.3.6]
\label{I.1.3.6}
For every $f\in A$, the open set $D(f)\subset X$ is canonically identified
with the prime spectrum $\operatorname{Spec}(A_f)$, and the sheaf
$\widetilde{M_f}$ associated with the $A_f$-module $M_f$ is canonically
identified with the restriction $\widetilde{M}|D(f)$.
\end{proposition}

The first assertion is a special case of
(\hyperref[I.1.2.6]{1.2.6}). Moreover, if $g\in A$ is such that
$D(g)\subset D(f)$, $M_g$ is canonically identified with the module of
fractions of $M_f$ whose denominators are the powers of the canonical image
of $g$ in $A_f$ (\hyperref[0.1.4.6]{0, 1.4.6}). The canonical
identification of $\widetilde{M_f}$ with $\widetilde{M}|D(f)$ then follows
from the definitions.

\begin{theorem}[1.3.7]
\label{I.1.3.7}
For every $A$-module $M$ and every $f\in A$, the homomorphism
\[
  \theta_f:M_f\longrightarrow\Gamma(D(f),\widetilde{M})
\]
is bijective (in other words, the presheaf $D(f)\mapsto M_f$ is a
\emph{sheaf}). In particular, $M$ is identified by means of $\theta_1$ with
$\Gamma(X,\widetilde{M})$.
\end{theorem}

Note that, if $M=A$, $\theta_f$ is a homomorphism of ring
structures; Theorem (\hyperref[I.1.3.7]{1.3.7}) will therefore imply that,
if we identify the rings $A_f$ and $\Gamma(D(f),\widetilde{A})$ by means of
$\theta_f$, the homomorphism
$\theta_f:M_f\to\Gamma(D(f),\widetilde{M})$ will be an isomorphism of
\emph{modules}.

Let us first show that $\theta_f$ is \emph{injective}. Indeed, if
$\xi\in M_f$ is such that $\theta_f(\xi)=0$, this means that for every prime
ideal $\mathfrak{p}$ of $A_f$, there exists $h\notin\mathfrak{p}$ such that
$h\xi=0$; since the annihilator of $\xi$ is contained in no prime ideal
of $A_f$, it is the whole of $A_f$, hence $\xi=0$.

It remains to show that $\theta_f$ is \emph{surjective}; we may reduce to
the case where $f=1$, the general case being deduced from it by ``localizing''
with the help of (\hyperref[I.1.3.6]{1.3.6}). So let $s$ be a section of
$\widetilde{M}$ over $X$; by virtue of (\hyperref[I.1.3.4]{1.3.4}) and of
(\hyperref[I.1.1.10]{1.1.10, (ii)}), there exists a \emph{finite}
covering $(D(f_i))_{i\in I}$ of $X$ ($f_i\in A$) such that, for every
$i\in I$, the restriction $s_i=s|D(f_i)$ is of the form
$\theta_{f_i}(\xi_i)$, where $\xi_i\in M_{f_i}$. If $i$, $j$ are two indices
in $I$, writing that the restrictions of $s_i$ and of $s_j$ to
$D(f_i)\cap D(f_j)=D(f_i f_j)$ are equal, we obtain, by definition of
$\widetilde{M}$,
\begin{equation}
  \rho_{f_i f_j,f_i}(\xi_i)=\rho_{f_i f_j,f_j}(\xi_j).
  \tag{1.3.7.1}\label{I.1.3.7.1}
\end{equation}
By definition, we can write, for each $i\in I$,
$\xi_i=z_i/f_i^{n_i}$, where $z_i\in M$, and since $I$ is finite, by
multiplying each $z_i$ by a power of $f_i$, we may assume all the $n_i$
equal to one and the same $n$. Then, by definition,
(\hyperref[I.1.3.7.1]{1.3.7.1}) means that there exists an integer
$m_{ij}\geq0$ such that $(f_i f_j)^{m_{ij}}(f_j^n z_i-f_i^n z_j)=0$, and we
may further assume all the $m_{ij}$ equal to one and the same integer $m$;
replacing $z_i$ by $f_i^m z_i$, we are then reduced to the case where $m=0$,
in other words to the case where we have
\begin{equation}
  f_j^n z_i=f_i^n z_j
  \tag{1.3.7.2}\label{I.1.3.7.2}
\end{equation}
for all $i$, $j$. Now, we have $D(f_i^n)=D(f_i)$, and since the $D(f_i)$
form a covering of $X$, the ideal generated by the $f_i^n$ is $A$;
in other words, there exist elements $g_i\in A$ such that
$\sum_i g_i f_i^n=1$. Consider then the element
$z=\sum_i g_i z_i$ of $M$; by (\hyperref[I.1.3.7.2]{1.3.7.2}), we have
$f_i^n z=\sum_j g_j f_i^n z_j=(\sum_j g_j f_j^n)z_i=z_i$, whence by
definition $\xi_i=z/1$ in $M_{f_i}$. We conclude
\oldpage[I]{87}
that $s_i$ is the restriction to $D(f_i)$ of $\theta_1(z)$, which proves
that $s=\theta_1(z)$ and completes the proof.

\begin{corollary}[1.3.8]
\label{I.1.3.8}
Let $M$, $N$ be two $A$-modules; the canonical homomorphism
$u\mapsto\widetilde{u}$ of $\operatorname{Hom}_A(M,N)$ into
$\operatorname{Hom}_{\widetilde{A}}(\widetilde{M},\widetilde{N})$ is
bijective. In particular, the relations $M=0$ and $\widetilde{M}=0$ are
equivalent.
\end{corollary}

Indeed, consider the canonical homomorphism $v\mapsto\Gamma(v)$ of
$\operatorname{Hom}_{\widetilde{A}}(\widetilde{M},\widetilde{N})$ into
$\operatorname{Hom}_{\Gamma(\widetilde{A})}
(\Gamma(\widetilde{M}),\Gamma(\widetilde{N}))$; the latter module is
canonically identified with $\operatorname{Hom}_A(M,N)$ by virtue of Theorem
(\hyperref[I.1.3.7]{1.3.7}). It remains to verify that
$u\mapsto\widetilde{u}$ and $v\mapsto\Gamma(v)$ are inverse to each
other; now, it is obvious that $\Gamma(\widetilde{u})=u$ by definition of
$\widetilde{u}$; on the other hand, if we set $u=\Gamma(v)$ for
$v\in\operatorname{Hom}_{\widetilde{A}}(\widetilde{M},\widetilde{N})$,
the map
$w:\Gamma(D(f),\widetilde{M})\to\Gamma(D(f),\widetilde{N})$ deduced
canonically from $v$ is such that the diagram
\[
  \xymatrix{
    M\ar[r]^u\ar[d]_{\rho_{f,1}} & N\ar[d]^{\rho_{f,1}}\\
    M_f\ar[r]_w & N_f
  }
\]
is commutative; we therefore necessarily have $w=u_f$ for every $f\in A$
(\hyperref[0.1.2.4]{0, 1.2.4}), which shows that
$\widetilde{\Gamma(v)}=v$.

\begin{corollary}[1.3.9]
\label{I.1.3.9}
\begin{enumerate}
  \item[\rm(i)] Let $u$ be a homomorphism of an $A$-module $M$ into an
  $A$-module $N$; then the sheaves associated with $\operatorname{Ker}u$,
  $\operatorname{Im}u$, $\operatorname{Coker}u$ are respectively
  $\operatorname{Ker}\widetilde{u}$, $\operatorname{Im}\widetilde{u}$,
  $\operatorname{Coker}\widetilde{u}$. In particular, for
  $\widetilde{u}$ to be injective (resp. surjective, bijective), it is
  necessary and sufficient that $u$ be so.
  \item[\rm(ii)] If $M$ is an inductive limit (resp. direct sum) of a
  family of $A$-modules $(M_\lambda)$, $\widetilde{M}$ is the inductive limit
  (resp. direct sum) of the family $(\widetilde{M_\lambda})$, up to a
  canonical isomorphism.
\end{enumerate}
\end{corollary}

(i) It suffices to apply the fact that $\widetilde{M}$ is an exact functor in
$M$ (\hyperref[I.1.3.5]{1.3.5}) to the two exact sequences of $A$-modules
\[
  0\to\operatorname{Ker}u\to M\to\operatorname{Im}u\to0
\]
\[
  0\to\operatorname{Im}u\to N\to\operatorname{Coker}u\to0
\]
The second assertion then follows from Theorem
(\hyperref[I.1.3.7]{1.3.7}).

(ii) Let $(M_\lambda,g_{\mu\lambda})$ be an inductive system of $A$-modules,
with inductive limit $M$, and let $g_\lambda$ be the canonical homomorphism
$M_\lambda\to M$. Since we have
$\widetilde{g_{\nu\mu}}\circ\widetilde{g_{\mu\lambda}}
=\widetilde{g_{\nu\lambda}}$ and
$\widetilde{g_\lambda}=\widetilde{g_\mu}\circ
\widetilde{g_{\mu\lambda}}$ for $\lambda\leq\mu\leq\nu$,
$(\widetilde{M_\lambda},\widetilde{g_{\mu\lambda}})$ is an inductive
system of sheaves on $X$, and if we denote by $h_\lambda$
the canonical homomorphism
$\widetilde{M_\lambda}\to\varinjlim\widetilde{M_\lambda}$, there is a
unique homomorphism
$v:\varinjlim\widetilde{M_\lambda}\to\widetilde{M}$ such that
$v\circ h_\lambda=\widetilde{g_\lambda}$. To see that $v$ is bijective, it
suffices to verify that, for every $x\in X$, $v_x$ is a bijection of
$(\varinjlim\widetilde{M_\lambda})_x$ onto $\widetilde{M}_x$; but
$\widetilde{M}_x=M_x$ and
\[
  (\varinjlim\widetilde{M_\lambda})_x
  =\varinjlim(\widetilde{M_\lambda})_x
  =\varinjlim(M_\lambda)_x=M_x
  \quad(\hyperref[0.1.3.3]{0, 1.3.3}).
\]
On the other hand, it follows from the definitions that
$(\widetilde{g_\lambda})_x$ and $(h_\lambda)_x$ are both equal to the
canonical map of $(M_\lambda)_x$ into $M_x$; since
$(\widetilde{g_\lambda})_x=v_x\circ(h_\lambda)_x$, $v_x$ is the identity.
\oldpage[I]{88}
Finally, if $M$ is the direct sum of two $A$-modules $N$, $P$, it is
immediate that $\widetilde{M}=\widetilde{N}\oplus\widetilde{P}$; since every
direct sum is the inductive limit of finite direct sums, the assertions of
(ii) are proved.

Note that (\hyperref[I.1.3.8]{1.3.8}) proves that the sheaves
isomorphic to sheaves associated with $A$-modules form an
\emph{abelian category} (T, I, 1.4).

Note also that it follows from (\hyperref[I.1.3.9]{1.3.9}) that if $M$
is an $A$-module \emph{of finite type}, that is to say, if there exists a
surjective homomorphism $A^n\to M$, then there exists a surjective
homomorphism $\widetilde{A^n}\to\widetilde{M}$; in other words, the
$\widetilde{A}$-module $\widetilde{M}$ is \emph{generated by a finite family
of sections over $X$} (\hyperref[0.5.1.1]{0, 5.1.1}), and conversely.

\begin{env}[1.3.10]
\label{I.1.3.10}
If $N$ is a submodule of an $A$-module $M$, the canonical injection
$j:N\to M$ gives, by (\hyperref[I.1.3.9]{1.3.9}), an injective
homomorphism $\widetilde{N}\to\widetilde{M}$, which makes it possible to
identify $\widetilde{N}$ canonically with a \emph{sub-$\widetilde{A}$-module}
of $\widetilde{M}$; we shall always assume this identification made. If $N$
and $P$ are two submodules of $M$, we then have
\begin{equation}
  \widetilde{(N+P)}=\widetilde{N}+\widetilde{P}
  \tag{1.3.10.1}\label{I.1.3.10.1}
\end{equation}
\begin{equation}
  \widetilde{(N\cap P)}=\widetilde{N}\cap\widetilde{P}
  \tag{1.3.10.2}\label{I.1.3.10.2}
\end{equation}
since $N+P$ and $N\cap P$ are respectively the image of the canonical
homomorphism $N\oplus P\to M$ and the kernel of the canonical homomorphism
$M\to(M/N)\oplus(M/P)$, and it suffices to apply
(\hyperref[I.1.3.9]{1.3.9}).

We conclude from (\hyperref[I.1.3.10.1]{1.3.10.1}) and
(\hyperref[I.1.3.10.2]{1.3.10.2}) that if
$\widetilde{N}=\widetilde{P}$, then $N=P$.
\end{env}

\begin{corollary}[1.3.11]
\label{I.1.3.11}
On the category of sheaves isomorphic to sheaves associated with
$A$-modules, the functor $\Gamma$ is exact.
\end{corollary}

Indeed, let
$\widetilde{M}\xrightarrow{\widetilde{u}}\widetilde{N}
\xrightarrow{\widetilde{v}}\widetilde{P}$ be an exact sequence corresponding
to two homomorphisms $u:M\to N$, $v:N\to P$ of $A$-modules. If
$Q=\operatorname{Im}u$ and $R=\operatorname{Ker}v$, we have
$\widetilde{Q}=\operatorname{Im}\widetilde{u}
=\operatorname{Ker}\widetilde{v}=\widetilde{R}$
(Corollary (\hyperref[I.1.3.9]{1.3.9})), hence $Q=R$.

\begin{corollary}[1.3.12]
\label{I.1.3.12}
Let $M$, $N$ be two $A$-modules.
\begin{enumerate}
  \item[\rm(i)] The sheaf associated with $M\otimes_A N$ is canonically
  identified with $\widetilde{M}\otimes_{\widetilde{A}}\widetilde{N}$.
  \item[\rm(ii)] If moreover $M$ admits a finite presentation, the sheaf
  associated with $\operatorname{Hom}_A(M,N)$ is canonically identified with
  $\mathscr{H}\!om_{\widetilde{A}}(\widetilde{M},\widetilde{N})$.
\end{enumerate}
\end{corollary}

(i) The sheaf
$\mathscr{F}=\widetilde{M}\otimes_{\widetilde{A}}\widetilde{N}$ is
associated with the presheaf
\[
  U\mapsto\mathscr{F}(U)=\Gamma(U,\widetilde{M})
  \otimes_{\Gamma(U,\widetilde{A})}\Gamma(U,\widetilde{N})
\]
$U$ running through the base (\hyperref[I.1.1.10]{1.1.10, (i)}) of $X$
consisting of the $D(f)$, where $f\in A$. Now, $\mathscr{F}(D(f))$ is
canonically identified with $M_f\otimes_{A_f}N_f$ by virtue of
(\hyperref[I.1.3.7]{1.3.7}) and (\hyperref[I.1.3.6]{1.3.6}). It is
known, moreover, that the $A_f$-module $M_f\otimes_{A_f}N_f$ is canonically
isomorphic to $(M\otimes_A N)_f$ (\hyperref[0.1.3.4]{0, 1.3.4}), which is
itself canonically isomorphic to
$\Gamma(D(f),\widetilde{(M\otimes_A N)})$
(\hyperref[I.1.3.7]{1.3.7} and \hyperref[I.1.3.6]{1.3.6}). Further,
one verifies at once that the canonical isomorphisms
\[
  \mathscr{F}(D(f))\simeq\Gamma(D(f),\widetilde{(M\otimes_A N)})
\]
\oldpage[I]{89}
thus obtained satisfy the conditions of compatibility with the restriction
operators (\hyperref[0.1.4.2]{0, 1.4.2}), and hence define a canonical
functorial isomorphism
\[
  \widetilde{M}\otimes_{\widetilde{A}}\widetilde{N}
  \simeq\widetilde{(M\otimes_A N)}
\]

(ii) The sheaf
$\mathscr{G}=\mathscr{H}\!om_{\widetilde{A}}(\widetilde{M},\widetilde{N})$
is associated with the presheaf
\[
  U\longmapsto\mathscr{G}(U)
  =\operatorname{Hom}_{\widetilde{A}|U}
  (\widetilde{M}|U,\widetilde{N}|U)
\]
$U$ running through the base of $X$ consisting of the $D(f)$. Now,
$\mathscr{G}(D(f))$ is canonically identified with
$\operatorname{Hom}_{A_f}(M_f,N_f)$
(\hyperref[I.1.3.6]{1.3.6} and \hyperref[I.1.3.8]{1.3.8}), which
is itself, by virtue of the hypothesis on $M$, canonically identified with
$(\operatorname{Hom}_A(M,N))_f$
(\hyperref[0.1.3.5]{0, 1.3.5}). Finally,
$(\operatorname{Hom}_A(M,N))_f$ is canonically identified with
$\Gamma(D(f),\widetilde{(\operatorname{Hom}_A(M,N))})$
(\hyperref[I.1.3.6]{1.3.6} and \hyperref[I.1.3.7]{1.3.7}), and the
canonical isomorphisms
$\mathscr{G}(D(f))\simeq
\Gamma(D(f),\widetilde{(\operatorname{Hom}_A(M,N))})$ thus obtained are
compatible with the restriction operators
(\hyperref[0.1.4.2]{0, 1.4.2}); they therefore define a canonical
isomorphism
\[
  \mathscr{H}\!om_{\widetilde{A}}(\widetilde{M},\widetilde{N})
  \simeq\widetilde{(\operatorname{Hom}_A(M,N))}.
\]

\begin{env}[1.3.13]
\label{I.1.3.13}
Now let $B$ be a (commutative) $A$-algebra; this can be interpreted by
saying that $B$ is an $A$-module and that we are given an element $e\in B$
and an $A$-homomorphism
$\varphi:B\otimes_A B\to B$ such that: $1^\circ$ the diagrams
\[
  \xymatrix{
    B\otimes_A B\otimes_A B\ar[r]^{\varphi\otimes 1}
      \ar[d]_{1\otimes\varphi} &
    B\otimes_A B\ar[d]^\varphi & &
    B\otimes_A B\ar[rr]^\sigma\ar[rd]_\varphi & &
    B\otimes_A B\ar[dl]^\varphi\\
    B\otimes_A B\ar[r]_\varphi &
    B & & & B
  }
\]
($\sigma$ being the canonical symmetry) are commutative; $2^\circ$
$\varphi(e\otimes x)=\varphi(x\otimes e)=x$. By virtue of
(\hyperref[I.1.3.12]{1.3.12}), the homomorphism
$\widetilde{\varphi}:\widetilde{B}\otimes_{\widetilde{A}}\widetilde{B}
\to\widetilde{B}$ of $\widetilde{A}$-modules satisfies the analogous
conditions, and therefore defines on $\widetilde{B}$ a structure of
$\widetilde{A}$-algebra. In the same way, giving a $B$-module $N$
amounts to giving an $A$-module $N$ and an $A$-homomorphism
$\psi:B\otimes_A N\to N$ such that the diagram
\[
  \xymatrix{
    B\otimes_A B\otimes_A N\ar[r]^{\varphi\otimes 1}
      \ar[d]_{1\otimes\psi} &
    B\otimes_A N\ar[d]^\psi\\
    B\otimes_A N\ar[r]_\psi & N
  }
\]
is commutative and $\psi(e\otimes n)=n$; the homomorphism
$\widetilde{\psi}:\widetilde{B}\otimes_{\widetilde{A}}\widetilde{N}
\to\widetilde{N}$ will satisfy the analogous condition, and therefore
defines on $\widetilde{N}$ a structure of $\widetilde{B}$-module.

In the same way, one sees that if $u:B\to B'$ (resp. $v:N\to N'$) is a
homomorphism of $A$-algebras (resp. of $B$-modules), $\widetilde{u}$
(resp. $\widetilde{v}$) is a homomorphism of $\widetilde{A}$-algebras
(resp. of $\widetilde{B}$-modules) and $\operatorname{Ker}\widetilde{u}$ is
a $\widetilde{B}$-Ideal (resp. $\operatorname{Ker}\widetilde{v}$,
$\operatorname{Coker}\widetilde{v}$ and $\operatorname{Im}\widetilde{v}$
are $\widetilde{B}$-modules). If $N$ is a $B$-module, $\widetilde{N}$ is
a $\widetilde{B}$-module of finite type if and only if $N$ is a
$B$-module of finite type (\hyperref[0.5.2.3]{0, 5.2.3}).
\oldpage[I]{90}
If $M$, $N$ are two $B$-modules, the $\widetilde{B}$-module
$\widetilde{M}\otimes_{\widetilde{B}}\widetilde{N}$ is canonically
identified with $\widetilde{(M\otimes_B N)}$; similarly,
$\mathscr{H}\!om_{\widetilde{B}}(\widetilde{M},\widetilde{N})$ is
canonically identified with $\widetilde{(\operatorname{Hom}_B(M,N))}$ when
$M$ admits a finite presentation; the proofs are the same as in
(\hyperref[I.1.3.12]{1.3.12}).

If $\mathfrak{J}$ is an ideal of $B$ and $N$ a $B$-module, we have
$\widetilde{(\mathfrak{J}N)}=\widetilde{\mathfrak{J}}\cdot\widetilde{N}$.

Finally, if $B$ is an $A$-algebra graded by sub-$A$-modules $B_n$
($n\in\mathbb{Z}$), the $\widetilde{A}$-algebra $\widetilde{B}$, the direct
sum of the $\widetilde{A}$-modules $\widetilde{B_n}$
(\hyperref[I.1.3.9]{1.3.9}), is graded by these
sub-$\widetilde{A}$-modules, the axiom of the grading expressing that the
image of the homomorphism $B_m\otimes_A B_n\to B$ is contained in
$B_{m+n}$. Similarly, if $M$ is a $B$-module graded by submodules $M_n$,
$\widetilde{M}$ is a $\widetilde{B}$-module graded by the
$\widetilde{M_n}$.
\end{env}

\begin{env}[1.3.14]
\label{I.1.3.14}
If $B$ is an $A$-algebra and $M$ a submodule of $B$, the
sub-$\widetilde{A}$-algebra of $\widetilde{B}$ generated by $\widetilde{M}$
(\hyperref[0.4.1.3]{0, 4.1.3}) is the sub-$\widetilde{A}$-algebra
$\widetilde{C}$, where $C$ denotes the subalgebra of $B$ generated by
$M$. Indeed, $C$ is the sum of the submodules of $B$ that are the images of
the homomorphisms $\bigotimes_A^n M\to B$ ($n\geq0$), and it suffices to
apply (\hyperref[I.1.3.9]{1.3.9}) and (\hyperref[I.1.3.12]{1.3.12}).
\end{env}


\subsection{Quasi-coherent sheaves on a prime spectrum.}
\label{subsection:I.1.4}

\begin{theorem}[1.4.1]
\label{I.1.4.1}
Let $X$ be the prime spectrum of a ring $A$, $V$ a quasi-compact open subset
of $X$, and $\mathscr{F}$ an $(\mathscr{O}_X|V)$-module. The following
four conditions are equivalent:
\begin{enumerate}
  \item[\textit{a})] There exists an $A$-module $M$ such that $\mathscr{F}$
  is isomorphic to $\widetilde{M}|V$.
  \item[\textit{b})] There exists a finite open covering $(V_i)$ of $V$ by
  sets of the form $D(f_i)$ ($f_i\in A$) contained in $V$, such that,
  for every $i$, $\mathscr{F}|V_i$ is isomorphic to a sheaf of the form
  $\widetilde{M_i}$, where $M_i$ is an $A_{f_i}$-module.
  \item[\textit{c})] The sheaf $\mathscr{F}$ is quasi-coherent
  (\hyperref[0.5.1.3]{0, 5.1.3}).
  \item[\textit{d})] The following two properties hold:
  \begin{enumerate}
    \item[\textit{d}\,1)] For every $f\in A$ such that $D(f)\subset V$ and
    every section $s\in\Gamma(D(f),\mathscr{F})$, there exists an integer
    $n\geq0$ such that $f^n s$ extends to a section of $\mathscr{F}$ over $V$.
    \item[\textit{d}\,2)] For every $f\in A$ such that $D(f)\subset V$ and
    every section $t\in\Gamma(V,\mathscr{F})$ such that the restriction of
    $t$ to $D(f)$ is $0$, there exists an integer $n\geq0$ such that
    $f^n t=0$.
  \end{enumerate}
\end{enumerate}
\end{theorem}

(In writing conditions \textit{d}\,1) and \textit{d}\,2), we have tacitly
identified $A$ and $\Gamma(\widetilde{A})$ by virtue of
\hyperref[I.1.3.7]{1.3.7}.)

That \textit{a}) implies \textit{b}) is an immediate consequence of
(\hyperref[I.1.3.6]{1.3.6}) and of the fact that the $D(f_i)$ form a base
for the topology of $X$ (\hyperref[I.1.1.10]{1.1.10}). Since every
$A$-module is isomorphic to the cokernel of a homomorphism
$A^{(I)}\to A^{(J)}$, (\hyperref[I.1.3.9]{1.3.9}) proves that every sheaf
associated with an $A$-module is quasi-coherent; hence \textit{b}) implies
\textit{c}). Conversely, if $\mathscr{F}$ is quasi-coherent, every $x\in V$
has a neighbourhood of the form $D(f)\subset V$ such that $\mathscr{F}|D(f)$
is isomorphic to the cokernel of a homomorphism
$\widetilde{A_f}^{(I)}\to\widetilde{A_f}^{(J)}$, hence to the sheaf
associated with the module $\widetilde{N}$, the cokernel of the corresponding
homomorphism $A_f^{(I)}\to A_f^{(J)}$
(\hyperref[I.1.3.8]{1.3.8} and \hyperref[I.1.3.9]{1.3.9}); since $V$
is quasi-compact, it is clear that \textit{c}) implies \textit{b}).
\oldpage[I]{91}
To prove that \textit{b}) implies \textit{d}\,1) and \textit{d}\,2),
suppose first that $V=D(g)$ for some $g\in A$, and that $\mathscr{F}$ is
isomorphic to a sheaf $\widetilde{N}$ associated with an $A_g$-module $N$;
replacing $X$ by $V$ and $A$ by $A_g$
(\hyperref[I.1.3.6]{1.3.6}), we may reduce to the case where $g=1$. Then
$\Gamma(D(f),\widetilde{N})$ and $N_f$ are canonically identified
(\hyperref[I.1.3.6]{1.3.6} and \hyperref[I.1.3.7]{1.3.7}), so a
section $s\in\Gamma(D(f),\widetilde{N})$ is identified with an element of
the form $z/f^n$, where $z\in N$; the section $f^n s$ is identified with the
element $z/1$ of $N_f$, and is consequently the restriction to $D(f)$ of the
section of $\widetilde{N}$ over $X$ identified with the element $z\in N$;
whence \textit{d}\,1) in this case. Similarly, $t\in\Gamma(X,\widetilde{N})$
is identified with an element $z'\in N$, the restriction of $t$ to $D(f)$ is
identified with the image $z'/1$ of $z'$ in $N_f$, and to say that this image
is zero means that there exists $n\geq0$ such that $f^n z'=0$ in $N$, or,
what amounts to the same, $f^n t=0$.

To finish proving that \textit{b}) implies \textit{d}\,1) and
\textit{d}\,2), it will suffice to establish the following lemma:

\begin{lemma}[1.4.1.1]
\label{I.1.4.1.1}
Suppose that $V$ is a finite union of sets of the form $D(g_i)$, and
that each of the sheaves $\mathscr{F}|D(g_i)$,
$\mathscr{F}|(D(g_i)\cap D(g_j))=\mathscr{F}|D(g_i g_j)$ satisfies
\textit{d}\,1) and \textit{d}\,2); then $\mathscr{F}$ has the following two
properties:
\begin{enumerate}
  \item[\textit{d}'\,1)] For every $f\in A$ and every section
  $s\in\Gamma(D(f)\cap V,\mathscr{F})$, there exists an integer $n\geq0$
  such that $f^n s$ extends to a section of $\mathscr{F}$ over $V$.
  \item[\textit{d}'\,2)] For every $f\in A$ and every section
  $t\in\Gamma(V,\mathscr{F})$ such that the restriction of $t$ to
  $D(f)\cap V$ is $0$, there exists an integer $n\geq0$ such that $f^n t=0$.
\end{enumerate}
\end{lemma}

Let us first prove \textit{d}'\,2): since
$D(f)\cap D(g_i)=D(fg_i)$, there exists for each $i$ an integer $n_i$ such
that the restriction of $(fg_i)^{n_i}t$ to $D(g_i)$ is zero: since
the image of $g_i$ in $A_{g_i}$ is invertible, the restriction of
$f^{n_i}t$ to $D(g_i)$ is also zero; taking for $n$ the largest of the
$n_i$, we have \textit{d}'\,2).

To prove \textit{d}'\,1), let us apply \textit{d}\,1) to the sheaf
$\mathscr{F}|D(g_i)$: there exist an integer $n_i\geq0$ and a section $s_i'$
of $\mathscr{F}$ over $D(g_i)$ extending the restriction of
$(fg_i)^{n_i}s$ to $D(fg_i)$; since the image of $g_i$ in $A_{g_i}$ is
invertible, there is a section $s_i$ of $\mathscr{F}$ over $D(g_i)$ such
that $s_i'=g_i^{n_i}s_i$, and $s_i$ extends the restriction of $f^{n_i}s$ to
$D(fg_i)$; we may moreover assume all the $n_i$ equal to one and the same
integer $n$. By construction, the restriction of $s_i-s_j$ to
$D(f)\cap D(g_i)\cap D(g_j)=D(fg_i g_j)$ is zero; by
\textit{d}\,2) applied to the sheaf $\mathscr{F}|D(g_i g_j)$, there exists an
integer $m_{ij}\geq0$ such that the restriction to $D(g_i g_j)$ of
$(fg_i g_j)^{m_{ij}}(s_i-s_j)$ is zero; since the image of $g_i g_j$
in $A_{g_i g_j}$ is invertible, the restriction of
$f^{m_{ij}}(s_i-s_j)$ to $D(g_i g_j)$ is zero. We may then assume all
the $m_{ij}$ equal to one and the same integer $m$, and there therefore
exists a section $s'\in\Gamma(V,\mathscr{F})$ extending the $f^m s_i$; this
section consequently extends $f^{n+m}s$, whence \textit{d}'\,1).

It remains to prove that \textit{d}\,1) and \textit{d}\,2) imply
\textit{a}). Let us first show that \textit{d}\,1) and \textit{d}\,2)
imply that these conditions are satisfied by every sheaf
$\mathscr{F}|D(g)$, where $g\in A$ is such that $D(g)\subset V$. This is
obvious for \textit{d}\,1); on the other hand, if $t\in\Gamma(D(g),\mathscr{F})$
is such that its restriction to $D(f)\subset D(g)$ is zero, there exists by
\textit{d}\,1) an integer $m\geq0$ such that $g^m t$
\oldpage[I]{92}
extends to a section $s$ of $\mathscr{F}$ over $V$; applying
\textit{d}\,2), we see that there exists an integer $n\geq0$ such that
$f^n g^m t=0$, and since the image of $g$ in $A_g$ is invertible, $f^n t=0$.

That being so, since $V$ is quasi-compact, Lemma
(\hyperref[I.1.4.1.1]{1.4.1.1}) proves that conditions
\textit{d}'\,1) and \textit{d}'\,2) are satisfied. Consider then the
$A$-module $M=\Gamma(V,\mathscr{F})$, and let us define a homomorphism of
$\widetilde{A}$-modules $u:\widetilde{M}\to j_*(\mathscr{F})$, where $j$ is
the canonical injection $V\to X$. Since the $D(f)$ form a base for the
topology of $X$, it suffices, for each $f\in A$, to define a
homomorphism
$u_f:M_f\to\Gamma(D(f),j_*(\mathscr{F}))=\Gamma(D(f)\cap V,\mathscr{F})$,
satisfying the usual compatibility conditions
(\hyperref[0.3.2.5]{0, 3.2.5}). Since the canonical image of $f$ in
$A_f$ is invertible, the restriction homomorphism
$M=\Gamma(V,\mathscr{F})\to\Gamma(D(f)\cap V,\mathscr{F})$ factors as
$M\longrightarrow M_f\xrightarrow{u_f}\Gamma(D(f)\cap V,\mathscr{F})$
(\hyperref[0.1.2.4]{0, 1.2.4}), and the verification of the compatibility
conditions for $D(g)\subset D(f)$ is immediate. That being so, let us show
that condition \textit{d}'\,1) (resp. \textit{d}'\,2)) implies that each
of the $u_f$ is surjective (resp. injective), which will prove that $u$ is
\emph{bijective}, and consequently that $\mathscr{F}$ is the restriction to
$V$ of an $\widetilde{A}$-module isomorphic to $\widetilde{M}$. Now, if
$s\in\Gamma(D(f)\cap V,\mathscr{F})$, there exists by \textit{d}'\,1) an
integer $n\geq0$ such that $f^n s$ extends to a section $z\in M$; we then
have $u_f(z/f^n)=s$, hence $u_f$ is surjective. Similarly, if $z\in M$ is
such that $u_f(z/1)=0$, this means that the restriction to $D(f)\cap V$ of
the section $z$ is zero; by \textit{d}'\,2), there exists an integer
$n\geq0$ such that $f^n z=0$, whence $z/1=0$ in $M_f$, and $u_f$ is
therefore injective.

\par\hfill Q.E.D.\par

\begin{corollary}[1.4.2]
\label{I.1.4.2}
Every quasi-coherent sheaf on a quasi-compact open subset of $X$ is induced
by a quasi-coherent sheaf on $X$.
\end{corollary}

\begin{corollary}[1.4.3]
\label{I.1.4.3}
Every quasi-coherent $\mathscr{O}_X$-algebra on
$X=\operatorname{Spec}(A)$ is isomorphic to an $\mathscr{O}_X$-algebra of
the form $\widetilde{B}$, where $B$ is an algebra over $A$; every
quasi-coherent $\widetilde{B}$-module is isomorphic to a
$\widetilde{B}$-module of the form $\widetilde{N}$, where $N$ is a
$B$-module.
\end{corollary}

Indeed, a quasi-coherent $\mathscr{O}_X$-algebra is a quasi-coherent
$\mathscr{O}_X$-module, hence of the form $\widetilde{B}$, where
$B$ is an $A$-module; the fact that $B$ is an $A$-algebra follows from the
characterization of the structure of $\mathscr{O}_X$-algebra by means of
the homomorphism
$\widetilde{B}\otimes_{\widetilde{A}}\widetilde{B}\to\widetilde{B}$ of
$\widetilde{A}$-modules, together with
(\hyperref[I.1.3.12]{1.3.12}). If $\mathscr{G}$ is a quasi-coherent
$\widetilde{B}$-module, it suffices to show that $\mathscr{G}$ is also a
quasi-coherent $\widetilde{A}$-module in order to conclude afterwards in the
same way; since the question is local, we may, by restricting ourselves to
an open subset of $X$ of the form $D(f)$, assume that $\mathscr{G}$ is the
cokernel of a homomorphism
$\widetilde{B}^{(I)}\to\widetilde{B}^{(J)}$ of $\widetilde{B}$-modules (and
\emph{a fortiori} of $\widetilde{A}$-modules); the proposition then follows
from (\hyperref[I.1.3.8]{1.3.8}) and
(\hyperref[I.1.3.9]{1.3.9}).

\subsection{Coherent sheaves on a prime spectrum.}
\label{subsection:I.1.5}

\begin{theorem}[1.5.1]
\label{I.1.5.1}
Let $A$ be a \emph{Noetherian} ring, $X=\operatorname{Spec}(A)$ its prime
spectrum, $V$ an open subset of $X$, and $\mathscr{F}$ an
$(\mathscr{O}_X|V)$-module. The following conditions are equivalent:
\begin{enumerate}
  \item[\textit{a})] $\mathscr{F}$ is coherent.
  \item[\textit{b})] $\mathscr{F}$ is of finite type and quasi-coherent.
  \item[\textit{c})] There exists an $A$-module $M$ of finite type such that
  $\mathscr{F}$ is isomorphic to the sheaf $\widetilde{M}|V$.
\end{enumerate}
\end{theorem}

\oldpage[I]{93}
\textit{a}) trivially implies \textit{b}). To see that \textit{b})
implies \textit{c}), let us first note, since $V$ is quasi-compact
(\hyperref[0.2.2.3]{0, 2.2.3}), that $\mathscr{F}$ is isomorphic to a
sheaf $\widetilde{N}|V$, where $N$ is an $A$-module
(\hyperref[I.1.4.1]{1.4.1}). Now, we have
$N=\varinjlim M_\lambda$, where $M_\lambda$ runs through the set of
sub-$A$-modules of finite type of $N$, whence
(\hyperref[I.1.3.9]{1.3.9})
$\mathscr{F}=\widetilde{N}|V=\varinjlim\widetilde{M_\lambda}|V$; but since
$\mathscr{F}$ is of finite type, and $V$ quasi-compact, there exists an index
$\lambda$ such that $\mathscr{F}=\widetilde{M_\lambda}|V$
(\hyperref[0.5.2.3]{0, 5.2.3}).

Let us finally show that \textit{c}) implies \textit{a}). It is clear that
$\mathscr{F}$ is then of finite type
(\hyperref[I.1.3.6]{1.3.6} and \hyperref[I.1.3.9]{1.3.9}); moreover,
the question being local, we may restrict ourselves to the case where
$V=D(f)$, $f\in A$. Since $A_f$ is Noetherian, we finally see that
everything comes down to proving that the kernel of a homomorphism
$\widetilde{A^n}\to\widetilde{M}$, where $M$ is an $A$-module, is of finite
type. Now, such a homomorphism is of the form $\widetilde{u}$, where $u$ is
a homomorphism $A^n\to M$ (\hyperref[I.1.3.8]{1.3.8}), and if
$P=\operatorname{Ker}u$, we have
$\widetilde{P}=\operatorname{Ker}\widetilde{u}$
(\hyperref[I.1.3.9]{1.3.9}). Since $A$ is Noetherian, $P$ is of finite
type, which completes the proof.

\begin{corollary}[1.5.2]
\label{I.1.5.2}
Under the hypotheses of (\hyperref[I.1.5.1]{1.5.1}), the sheaf
$\mathscr{O}_X$ is a coherent sheaf of rings.
\end{corollary}

\begin{corollary}[1.5.3]
\label{I.1.5.3}
Under the hypotheses of (\hyperref[I.1.5.1]{1.5.1}), every coherent
sheaf on an open subset of $X$ is induced by a coherent sheaf on $X$.
\end{corollary}

\begin{corollary}[1.5.4]
\label{I.1.5.4}
Under the hypotheses of (\hyperref[I.1.5.1]{1.5.1}), every
quasi-coherent $\mathscr{O}_X$-module $\mathscr{F}$ is the inductive limit
of the coherent sub-$\mathscr{O}_X$-modules of $\mathscr{F}$.
\end{corollary}

Indeed, $\mathscr{F}=\widetilde{M}$, where $M$ is an $A$-module, and $M$ is
the inductive limit of its submodules of finite type; we conclude by
(\hyperref[I.1.3.9]{1.3.9}) and
(\hyperref[I.1.5.1]{1.5.1}).

\subsection{Functorial properties of quasi-coherent sheaves on a prime
spectrum.}
\label{subsection:I.1.6}

\begin{env}[1.6.1]
\label{I.1.6.1}
Let $A$, $A'$ be two rings,
\[
  \varphi:A'\to A
\]
a homomorphism, and
\[
  {}^a\varphi:X=\operatorname{Spec}(A)\to
  X'=\operatorname{Spec}(A')
\]
the continuous map associated with $\varphi$
(\hyperref[I.1.2.1]{1.2.1}). We are going to define a
\emph{canonical homomorphism}
\[
  \widetilde{\varphi}:\mathscr{O}_{X'}\to
  {}^a\varphi_*(\mathscr{O}_X)
\]
of sheaves of rings. For every $f'\in A'$, set $f=\varphi(f')$; we have
${}^a\varphi^{-1}(D(f'))=D(f)$
(\hyperref[I.1.2.2.2]{1.2.2.2}). The rings
$\Gamma(D(f'),\widetilde{A'})$ and $\Gamma(D(f),\widetilde{A})$
are identified with $A'_{f'}$ and $A_f$ respectively
(\hyperref[I.1.3.6]{1.3.6} and \hyperref[I.1.3.7]{1.3.7}). Now,
the homomorphism $\varphi$ canonically defines a homomorphism
$\varphi_{f'}:A'_{f'}\to A_f$
(\hyperref[0.1.5.1]{0, 1.5.1}); in other words, we have a ring
homomorphism
\[
  \Gamma(D(f'),\widetilde{A'})\to
  \Gamma({}^a\varphi^{-1}(D(f')),\widetilde{A})
  =\Gamma(D(f'),{}^a\varphi_*(\widetilde{A})).
\]

\oldpage[I]{94}
Moreover, these homomorphisms satisfy the usual compatibility
conditions: for $D(f')\supset D(g')$, the diagram
\[
  \xymatrix{
    \Gamma(D(f'),\widetilde{A'})\ar[r]\ar[d] &
    \Gamma(D(f'),{}^a\varphi_*(\widetilde{A}))\ar[d]\\
    \Gamma(D(g'),\widetilde{A'})\ar[r] &
    \Gamma(D(g'),{}^a\varphi_*(\widetilde{A}))
  }
\]
is commutative (\hyperref[0.1.5.1]{0, 1.5.1}); we have therefore indeed
defined a homomorphism of $\mathscr{O}_{X'}$-algebras, the $D(f')$ forming a
base for the topology of $X'$ (\hyperref[0.3.2.3]{0, 3.2.3}). The pair
$\Phi=({}^a\varphi,\widetilde{\varphi})$ is therefore a \emph{morphism}
of ringed spaces
\[
  \Phi:(X,\mathscr{O}_X)\to(X',\mathscr{O}_{X'})
\]
(\hyperref[0.4.1.1]{0, 4.1.1}).

Note moreover that, if we set $x'={}^a\varphi(x)$, the homomorphism
$\widetilde{\varphi}_x^\#$ (\hyperref[0.3.7.1]{0, 3.7.1}) is none other
than the homomorphism
\[
  \varphi_x:A'_{x'}\to A_x
\]
deduced canonically from $\varphi:A'\to A$
(\hyperref[0.1.5.1]{0, 1.5.1}). Indeed, every $z'\in A'_{x'}$
can be written $g'/f'$, where $f'$, $g'$ are in $A'$ and
$f'\notin\mathfrak{j}_{x'}$; $D(f')$ is therefore a neighbourhood of $x'$ in
$X'$, and the homomorphism
$\Gamma(D(f'),\widetilde{A'})\to
\Gamma({}^a\varphi^{-1}(D(f')),\widetilde{A})$ deduced from
$\widetilde{\varphi}$ is none other than $\varphi_{f'}$; by considering the
section $s'\in\Gamma(D(f'),\widetilde{A'})$ corresponding to
$g'/f'\in A'_{f'}$, we indeed obtain
$\widetilde{\varphi}_x^\#(z')=\varphi(g')/\varphi(f')$ in $A_x$.
\end{env}

\begin{example}[1.6.2]
\label{I.1.6.2}
Let $S$ be a multiplicative subset of $A$ and $\varphi$ the canonical
homomorphism $A\to S^{-1}A$; we have already seen
(\hyperref[I.1.2.6]{1.2.6}) that ${}^a\varphi$ is a
\emph{homeomorphism} of $Y=\operatorname{Spec}(S^{-1}A)$ onto the
subspace of $X=\operatorname{Spec}(A)$ consisting of the $x$ such that
$\mathfrak{j}_x\cap S=\varnothing$. Moreover, for every $x$ belonging to
this subspace, hence of the form ${}^a\varphi(y)$ with $y\in Y$,
the homomorphism
$\widetilde{\varphi}_y^\#:\mathscr{O}_x\to\mathscr{O}_y$ is
\emph{bijective} (\hyperref[0.1.2.6]{0, 1.2.6}); in other words,
$\mathscr{O}_Y$ is identified with the sheaf induced on $Y$ by
$\mathscr{O}_X$.
\end{example}

\begin{proposition}[1.6.3]
\label{I.1.6.3}
For every $A$-module $M$, there exists a canonical functorial isomorphism of
the $\mathscr{O}_{X'}$-module $(M_{[\varphi]})^{\sim}$ onto the direct image
$\Phi_*(\widetilde{M})$.
\end{proposition}

To abbreviate, set $M'=M_{[\varphi]}$, and for every $f'\in A'$, set
again $f=\varphi(f')$. The modules of sections
$\Gamma(D(f'),\widetilde{M'})$ and $\Gamma(D(f),\widetilde{M})$
are identified respectively with the modules $M'_{f'}$ and $M_f$ (over
$A'_{f'}$ and $A_f$ respectively); moreover, the $A'_{f'}$-module
$(M_f)_{[\varphi_{f'}]}$ is canonically isomorphic to $M'_{f'}$
(\hyperref[0.1.5.2]{0, 1.5.2}). We therefore have a functorial
isomorphism of $\Gamma(D(f'),\widetilde{A'})$-modules:
$\Gamma(D(f'),\widetilde{M'})\xrightarrow{\sim}
\Gamma({}^a\varphi^{-1}(D(f')),\widetilde{M})_{[\varphi_{f'}]}$, and these
isomorphisms satisfy the usual conditions of compatibility with the
restrictions (\hyperref[0.1.5.6]{0, 1.5.6}), hence define the announced
functorial isomorphism. Note that, more precisely, if
$u:M_1\to M_2$ is a homomorphism of $A$-modules, it can be regarded as a
homomorphism $(M_1)_{[\varphi]}\to(M_2)_{[\varphi]}$ of
$A'$-modules; if we denote this homomorphism by $u_{[\varphi]}$,
$\Phi_*(\widetilde{u})$ is identified with $(u_{[\varphi]})^{\sim}$.

This proof also shows that for every $A$-algebra $B$, the
\oldpage[I]{95}
canonical functorial isomorphism
$(B_{[\varphi]})^{\sim}\xrightarrow{\sim}\Phi_*(\widetilde{B})$ is an
isomorphism of $\mathscr{O}_{X'}$-algebras; if $M$ is a $B$-module,
the canonical functorial isomorphism
$(M_{[\varphi]})^{\sim}\xrightarrow{\sim}\Phi_*(\widetilde{M})$ is an
isomorphism of $\Phi_*(\widetilde{B})$-modules.

\begin{corollary}[1.6.4]
\label{I.1.6.4}
The direct image functor $\Phi_*$ is exact on the category of
quasi-coherent $\mathscr{O}_X$-modules.
\end{corollary}

Indeed, it is clear that $M_{[\varphi]}$ is an exact functor in $M$, and
$\widetilde{M'}$ an exact functor in $M'$ (\hyperref[I.1.3.5]{1.3.5}).

\begin{proposition}[1.6.5]
\label{I.1.6.5}
Let $N'$ be an $A'$-module and $N$ the $A$-module
$N'\otimes_{A'}A_{[\varphi]}$; there exists a canonical functorial
isomorphism of the $\mathscr{O}_X$-module $\Phi^*(\widetilde{N'})$ onto
$\widetilde{N}$.
\end{proposition}

Observe first that $j:z'\mapsto z'\otimes 1$ is an
$A'$-homomorphism of $N'$ into $N_{[\varphi]}$; indeed, by definition,
for $f'\in A'$, we have
$(f'z')\otimes 1=z'\otimes\varphi(f')=\varphi(f')(z'\otimes 1)$. We
deduce from it (\hyperref[I.1.3.8]{1.3.8}) a homomorphism
$\widetilde{j}:\widetilde{N'}\to(N_{[\varphi]})^{\sim}$ of
$\mathscr{O}_{X'}$-modules, and by virtue of
(\hyperref[I.1.6.3]{1.6.3}), we may regard $\widetilde{j}$ as
mapping $\widetilde{N'}$ into $\Phi_*(\widetilde{N})$. To this homomorphism
$\widetilde{j}$ there corresponds canonically a homomorphism
$h=\widetilde{j}^{\#}$ of $\Phi^*(\widetilde{N'})$ into $\widetilde{N}$
(\hyperref[0.4.4.3]{0, 4.4.3}); we shall see that on each stalk,
$h_x$ is \emph{bijective}. Set $x'={}^a\varphi(x)$ and let $f'\in A'$ be
such that $x'\in D(f')$; let $f=\varphi(f')$. The ring
$\Gamma(D(f),\widetilde{A})$ is identified with $A_f$, and the modules
$\Gamma(D(f),\widetilde{N})$ and $\Gamma(D(f'),\widetilde{N'})$ with $N_f$
and $N'_{f'}$ respectively; let
$s'\in\Gamma(D(f'),\widetilde{N'})$, identified with $n'/{f'}^p$
($n'\in N'$), and let $s$ be its image under $\widetilde{j}$ in
$\Gamma(D(f),\widetilde{N})$; $s$ is identified with $(n'\otimes 1)/f^p$.
On the other hand, let $t\in\Gamma(D(f),\widetilde{A})$, identified with
$g/f^q$ ($g\in A$); then, by definition, we have
$h_x(s'_{x'}\otimes t_x)=t_x\cdot s_x$
(\hyperref[0.4.4.3]{0, 4.4.3}). But we can canonically identify
$N_f$ with
$N'_{f'}\otimes_{A'_{f'}}(A_f)_{[\varphi_{f'}]}$
(\hyperref[0.1.5.4]{0, 1.5.4}); $s$ then corresponds to the element
$(n'/{f'}^p)\otimes 1$, and the section $y\mapsto t_y\cdot s_y$ to
$(n'/{f'}^p)\otimes(g/f^q)$. The compatibility diagrams of
(\hyperref[0.1.5.6]{0, 1.5.6}) show that $h_x$ is none other than
the canonical isomorphism
\[
  \label{I.1.6.5.1}
  N'_{x'}\otimes_{A'_{x'}}(A_x)_{[\varphi_x]}
  \xrightarrow{\sim}N_x=(N'\otimes_{A'}A_{[\varphi]})_x.
  \tag{1.6.5.1}
\]

Moreover, let $v:N'_1\to N'_2$ be a homomorphism of $A'$-modules; since
$\widetilde{v}_{x'}=v_{x'}$ for every $x'\in X'$, it follows at once from
what precedes that $\Phi^*(\widetilde{v})$ is canonically identified with
$(v\otimes 1)^{\sim}$, which completes the proof of
(\hyperref[I.1.6.5]{1.6.5}).

If $B'$ is an $A'$-algebra, the canonical isomorphism of
$\Phi^*(\widetilde{B'})$ onto
$(B'\otimes_{A'}A_{[\varphi]})^{\sim}$ is an isomorphism of
$\mathscr{O}_X$-algebras; if moreover $N'$ is a $B'$-module,
the canonical isomorphism of $\Phi^*(\widetilde{N'})$ onto
$(N'\otimes_{A'}A_{[\varphi]})^{\sim}$ is an isomorphism of
$\Phi^*(\widetilde{B'})$-modules.

\begin{corollary}[1.6.6]
\label{I.1.6.6}
The sections of $\Phi^*(\widetilde{N'})$ that are the canonical images of
the sections $s'$, where $s'$ runs through the $A'$-module
$\Gamma(\widetilde{N'})$, generate the $A$-module
$\Gamma(\Phi^*(\widetilde{N'}))$.
\end{corollary}

Indeed, these images are identified with the elements $z'\otimes 1$ of $N$
when $N'$ and $N$ are identified with $\Gamma(\widetilde{N'})$ and
$\Gamma(\widetilde{N})$ respectively (\hyperref[I.1.3.7]{1.3.7}) and
$z'$ runs through $N'$.

\begin{env}[1.6.7]
\label{I.1.6.7}
In the proof of (\hyperref[I.1.6.5]{1.6.5}), we proved in
passing that the canonical map
(\hyperref[0.4.4.3.2]{0, 4.4.3.2})
$\rho:\widetilde{N'}\to\Phi_*(\Phi^*(\widetilde{N'}))$ is none other than
the homomorphism $\widetilde{j}$,
\oldpage[I]{96}
where $j:N'\to N'\otimes_{A'}A_{[\varphi]}$ is the homomorphism
$z'\mapsto z'\otimes 1$. Similarly, the canonical map
(\hyperref[0.4.4.3.3]{0, 4.4.3.3})
$\sigma:\Phi^*(\Phi_*(\widetilde{M}))\to\widetilde{M}$ is none other than
$\widetilde{p}$, where
$p:M_{[\varphi]}\otimes_{A'}A_{[\varphi]}\to M$ is the canonical
homomorphism which maps every tensor product $z\otimes a$ ($z\in M$,
$a\in A$) to $a\cdot z$; this follows at once from the definitions
(\hyperref[0.3.7.1]{0, 3.7.1},
\hyperref[0.4.4.3]{0, 4.4.3}, and
\hyperref[I.1.3.7]{I, 1.3.7}).

We conclude from this (\hyperref[0.4.4.3]{0, 4.4.3} and
\hyperref[0.3.5.4.4]{3.5.4.4}) that if
$v:N'\to M_{[\varphi]}$ is an $A'$-homomorphism, we have
$\widetilde{v}^{\#}=(v\otimes 1)^{\sim}$.
\end{env}

\begin{env}[1.6.8]
\label{I.1.6.8}
Let $N'_1$, $N'_2$ be two $A'$-modules, $N'_1$ being assumed to have a
\emph{finite presentation}; it then follows from
(\hyperref[I.1.6.7]{1.6.7}) and
(\hyperref[I.1.3.12]{1.3.12, (ii)}) that the canonical homomorphism
(\hyperref[0.4.4.6]{0, 4.4.6})
\[
  \Phi^*\bigl(\mathscr{H}\!om_{\widetilde{A'}}
    (\widetilde{N'_1},\widetilde{N'_2})\bigr)
  \longrightarrow
  \mathscr{H}\!om_{\widetilde A}
    \bigl(\Phi^*(\widetilde{N'_1}),\Phi^*(\widetilde{N'_2})\bigr)
\]
is none other than $\widetilde{\gamma}$, where $\gamma$ denotes
the canonical homomorphism of $A$-modules
\[
  \operatorname{Hom}_{A'}(N'_1,N'_2)\otimes_{A'}A
  \longrightarrow
  \operatorname{Hom}_A(N'_1\otimes_{A'}A,N'_2\otimes_{A'}A).
\]
\end{env}

\begin{env}[1.6.9]
\label{I.1.6.9}
Let $\mathfrak{J}'$ be an ideal of $A'$ and $M$ an $A$-module; since, by
definition, $\widetilde{\mathfrak{J}'}\widetilde M$ is the image of
the canonical homomorphism
$\Phi^*(\widetilde{\mathfrak{J}'})\otimes_{\widetilde A}\widetilde M
\to\widetilde M$, it follows from (\hyperref[I.1.6.5]{1.6.5}) and
(\hyperref[I.1.3.12]{1.3.12, (i)}) that
$\widetilde{\mathfrak{J}'}\widetilde M$ is canonically identified with
$(\mathfrak{J}'M)^{\sim}$; in particular,
$\Phi^*(\widetilde{\mathfrak{J}'})\widetilde A$ is identified with
$(\mathfrak{J}'A)^{\sim}$, and, taking into account the right exactness of
the functor $\Phi^*$, the $\widetilde A$-algebra
$\Phi^*((A'/\mathfrak{J}')^{\sim})$ is identified with
$(A/\mathfrak{J}'A)^{\sim}$.
\end{env}

\begin{env}[1.6.10]
\label{I.1.6.10}
Let $A''$ be a third ring and $\varphi'$ a homomorphism
$A''\to A'$, and set $\varphi''=\varphi\circ\varphi'$. It follows
at once from the definitions that
${}^a\varphi''=({}^a\varphi')\circ({}^a\varphi)$ and
$\widetilde{\varphi''}=\widetilde\varphi\circ\widetilde{\varphi'}$
(\hyperref[0.1.5.7]{0, 1.5.7}). We conclude that we have
$\Phi''=\Phi'\circ\Phi$; in other words,
$(\operatorname{Spec} A,\widetilde A)$ is a \emph{functor} in $A$ from the
category of rings to that of ringed spaces.
\end{env}

\subsection{Characterization of morphisms of affine schemes.}
\label{subsection:I.1.7}

\begin{definition}[1.7.1]
\label{I.1.7.1}
A ringed space $(X,\mathscr{O}_X)$ is called an \emph{affine scheme} if it
is isomorphic to a ringed space of the form
$(\operatorname{Spec}(A),\widetilde A)$, where $A$ is a ring; we then say
that $\Gamma(X,\mathscr{O}_X)$, which is canonically identified with the ring
$A$ (\hyperref[I.1.3.7]{1.3.7}), is the ring of the affine scheme
$(X,\mathscr{O}_X)$, and we denote it by $A(X)$ when no confusion can
result.
\end{definition}

By abuse of language, when we speak of the \emph{affine scheme}
$\operatorname{Spec}(A)$, we shall always mean the ringed space
$(\operatorname{Spec}(A),\widetilde A)$.

\begin{env}[1.7.2]
\label{I.1.7.2}
Let $A$, $B$ be two rings, and $(X,\mathscr{O}_X)$,
$(Y,\mathscr{O}_Y)$ the corresponding affine schemes on the prime
spectra $X=\operatorname{Spec}(A)$, $Y=\operatorname{Spec}(B)$. We have
seen (\hyperref[I.1.6.1]{1.6.1}) that to every ring homomorphism
$\varphi:B\to A$ there corresponds a morphism
$\Phi=({}^a\varphi,\widetilde\varphi)=\operatorname{Spec}(\varphi):
(X,\mathscr{O}_X)\to(Y,\mathscr{O}_Y)$. Note that $\varphi$ is
entirely determined by $\Phi$, since by definition we have
\[
  \varphi=\Gamma(\widetilde\varphi):
  \Gamma(\widetilde B)\longrightarrow
  \Gamma({}^a\varphi_*(\widetilde A))=\Gamma(\widetilde A).
\]
\end{env}

\begin{theorem}[1.7.3]
\label{I.1.7.3}
Let $(X,\mathscr{O}_X)$, $(Y,\mathscr{O}_Y)$ be two affine schemes. For a
morphism of ringed spaces
$(\psi,\theta):(X,\mathscr{O}_X)\to(Y,\mathscr{O}_Y)$ to be of the form
$({}^a\varphi,\widetilde\varphi)$, where $\varphi$ is a ring homomorphism
$A(Y)\to A(X)$, it is necessary and sufficient that, for every $x\in X$,
$\theta_x^{\#}$ be a local homomorphism:
$\mathscr{O}_{\psi(x)}\to\mathscr{O}_x$.
\end{theorem}

\oldpage[I]{97}
Set $A=A(X)$, $B=A(Y)$. The condition is necessary, since we have seen
(\hyperref[I.1.6.1]{1.6.1}) that $\widetilde\varphi_x^{\#}$ is
the homomorphism of $B_{{}^a\varphi(x)}$ into $A_x$ deduced canonically from
$\varphi$, and, by definition of
${}^a\varphi(x)=\varphi^{-1}(\mathfrak{j}_x)$, this homomorphism is local.

Let us prove that the condition is sufficient. By definition, $\theta$ is a
homomorphism $\mathscr{O}_Y\to\psi_*(\mathscr{O}_X)$, and from it we deduce
canonically a ring homomorphism
\[
  \varphi=\Gamma(\theta):
  B=\Gamma(Y,\mathscr{O}_Y)
  \longrightarrow
  \Gamma(Y,\psi_*(\mathscr{O}_X))=\Gamma(X,\mathscr{O}_X)=A.
\]

The hypothesis on $\theta_x^{\#}$ makes it possible to deduce from this
homomorphism, by passing to the quotients, a monomorphism $\theta^x$ of the
residue field $k(\psi(x))$ into the residue field $k(x)$ such that, for every
section $f\in\Gamma(Y,\mathscr{O}_Y)=B$, we have
$\theta^x(f(\psi(x)))=\varphi(f)(x)$. The relation $f(\psi(x))=0$ is
therefore equivalent to $\varphi(f)(x)=0$, which means that
$\mathfrak{j}_{\psi(x)}=\mathfrak{j}_{{}^a\varphi(x)}$, which can also be
written $\psi(x)={}^a\varphi(x)$ for every $x\in X$, or $\psi={}^a\varphi$.
We also know that the diagram
\phantomsection\label{I.1.7.3.diagram}
\[
  \xymatrix{
    B=\Gamma(Y,\mathscr{O}_Y)\ar[r]^\varphi\ar[d] &
    \Gamma(X,\mathscr{O}_X)=A\ar[d]\\
    B_{\psi(x)}\ar[r]_{\theta_x^{\#}} & A_x
  }
\]
is commutative (\hyperref[0.3.7.2]{0, 3.7.2}), which means that
$\theta_x^{\#}$ is equal to the homomorphism
$\varphi_x:B_{\psi(x)}\to A_x$ deduced canonically from $\varphi$
(\hyperref[0.1.5.1]{0, 1.5.1}). Since the data of the $\theta_x^{\#}$
completely characterize $\theta^{\#}$, and consequently also $\theta$
(\hyperref[0.3.7.1]{0, 3.7.1}), we conclude that
$\theta=\widetilde\varphi$, by definition of $\widetilde\varphi$
(\hyperref[I.1.6.1]{1.6.1}).

We shall say that a morphism $(\psi,\theta)$ of ringed spaces satisfying
the condition of (\hyperref[I.1.7.3]{1.7.3}) is a \emph{morphism of
affine schemes}.

\begin{corollary}[1.7.4]
\label{I.1.7.4}
If $(X,\mathscr{O}_X)$, $(Y,\mathscr{O}_Y)$ are two affine schemes, there
exists a canonical isomorphism of the set of morphisms of affine
schemes
$\operatorname{Hom}((X,\mathscr{O}_X),(Y,\mathscr{O}_Y))$ onto the set
of ring homomorphisms of $B$ into $A$, where
$A=\Gamma(\mathscr{O}_X)$ and $B=\Gamma(\mathscr{O}_Y)$.
\end{corollary}

One can also say that the functors $(\operatorname{Spec}(A),\widetilde A)$
in $A$ and $\Gamma(X,\mathscr{O}_X)$ in $(X,\mathscr{O}_X)$ define an
\emph{equivalence} between the category of commutative rings and the
category dual to the category of affine schemes (T, I, 1.2).

\begin{corollary}[1.7.5]
\label{I.1.7.5}
If $\varphi:B\to A$ is surjective, the corresponding morphism
$({}^a\varphi,\widetilde\varphi)$ is a monomorphism of ringed spaces
\emph{(cf. (\hyperref[I.4.1.7]{4.1.7}))}.
\end{corollary}

Indeed, we know that ${}^a\varphi$ is injective
(\hyperref[I.1.2.5]{1.2.5}), and since $\varphi$ is surjective, for every
$x\in X$, $\varphi_x^{\#}:B_{{}^a\varphi(x)}\to A_x$, which is deduced from
$\varphi$ by passing to the rings of fractions, is also surjective
(\hyperref[0.1.5.1]{0, 1.5.1}); whence the conclusion
(\hyperref[0.4.1.1]{0, 4.1.1}).
